\section{Méthodologie de test et de validation}

La fiabilité du système repose autant sur la méthode de vérification employée tout au long du projet que sur les correctifs eux-mêmes. Trois principes ont été appliqués de façon constante.

\textbf{L'inspection directe des données prévaut toujours sur la supposition.} Avant de conclure qu'une réponse est correcte ou incorrecte, le contenu réel de Neo4j et de Qdrant a systématiquement été interrogé pour établir une vérité de référence -- y compris lorsque cela signifiait reproduire fidèlement une coquille présente dans les données sources plutôt que de la ``corriger'' silencieusement dans la réponse générée. Un audit exhaustif de complétude a par ailleurs été mené en comparant, entité par entité, le contenu des deux bases aux fichiers sources d'origine, ce qui a directement mené à la découverte de l'incident de couverture Centre-Service décrit en section~\ref{sec:incidents}.

\textbf{Chaque correctif est vérifié par un nouveau test réel, jamais supposé résolu sur la seule lecture du code.} Après chaque modification, la pile complète est redémarrée (arrêt puis démarrage, jamais un redémarrage partiel) et le cache Redis est systématiquement vidé avant de rejouer le scénario ayant révélé le défaut, pour exclure qu'une ancienne réponse mise en cache ne masque une régression ou ne simule à tort une correction.

\textbf{Les tests couvrent délibérément la diversité réelle des formulations citoyennes.} Au-delà des scénarios de démonstration, un jeu de quinze questions a été construit pour représenter la variété réelle des façons de poser une question~: registre formel et darija familière, arabe et français, questions directes et questions elliptiques ou vagues, questions à critères multiples, question de suivi conversationnel, question hors domaine, calcul arithmétique, et un cas de négation (limitation documentée ci-dessous). Sur ce jeu, treize questions sur quinze ont abouti à une réponse jugée correcte et utile après comparaison avec le contenu réel de la base ; les défauts identifiés (réponse de repli inexploitable pour les FAQ/programmes, une formulation familière de demande de crèche mal reconnue, la limitation de négation, un dépassement de délai sur une question à critères multiples) ont chacun été soit corrigés et re-vérifiés (cas de la réponse de repli, section~\ref{sec:incidents}), soit documentés comme limitation connue plutôt que dissimulés.
